\chapter{Introduction}
\label{chap:intro}

\section{Context and Motivation}
\label{sec:intro-context}

Online retail has outgrown what any one shopper can browse. A
single category page on a large marketplace lists thousands of
substitutes, and shoppers spend a lot of time browsing, comparing,
and validating products before they buy. Recommender systems
narrow the search space using signals like co-purchase frequency,
content similarity over product attributes, and ranking models
trained on historical
interactions~\cite{ComprehensiveRecSys2024, LLM4Rec2025}. These
methods are well studied, fast, and widely deployed. They share
one limitation that hits the user directly: the output is opaque.
A shopper who sees a ``frequently bought together'' or
``similar items'' panel is not told why those products were
chosen.

Large language models reopen the explanation problem. An LLM can
write natural-language rationales, summarize customer reviews,
and surface tradeoffs between similar
products~\cite{LLMEnhancedRecSys2024, ExplainableRecKG2024}.
Putting an LLM into a shopping flow brings new risks, though.
Amazon Rufus is the most visible case in production. The system
serves over 300 million users and is credited with roughly
\$12 billion in incremental sales~\cite{RufusRevenue2025}.
Despite that scale, independent studies report a 32\% accuracy
rate for picking the right product, a 28\% price hallucination
rate, and an 83\% bias toward Amazon-owned
products~\cite{RufusWrong2024,RufusConfidentlyWrong2024}.
Consumer surveys back this up. 55\% of users say they do not
trust AI chatbot product
recommendations~\cite{ConsumerDistrust2025}, and 42\% see
e-commerce AI assistants as upsell tools, not
advisors~\cite{AIUpsell2025}. HCI research adds a timing angle:
AI suggestions offered after a natural pause land better than
proactive interruptions~\cite{CHI2024Timing}.

The open question is not whether to put LLMs into e-commerce,
but how. The architecture and interaction choices that decide
whether the assistant is trusted, accurate, and welcome have not
been studied systematically at bachelor-thesis scale against a
documented production baseline.

\section{Motivation}
\label{sec:intro-motivation}

Two facts shape the design space. First, the failure modes of
production AI shopping assistants are documented, not
speculative. Hallucinated product data, fabricated review themes,
and perceived intrusiveness all show up in Rufus and in user
surveys, and they are not edge
cases~\cite{RufusWrong2024,ConsumerResistanceAI}. Second, the
timing of an AI suggestion matters. Suggestions offered after a
natural pause read as helpful. Suggestions injected during active
deliberation read as
interruption~\cite{CHI2024Timing,ProactiveReactivePersonalization}.

A reactive-first interaction model follows. The system tracks
lightweight browsing signals (dwell time on a product, scroll
depth, engagement with the reviews tab) to infer when help might
be welcome, but it never auto-fires AI advice. The user has to
engage on purpose. Behavior signals prepare context. The user
pulls.

Amazon's own product timeline backs this up. Rufus launched as a
reactive assistant. The proactive ``Help Me Decide'' surface was
added ten months later, after the reactive product had
stabilized~\cite{RufusTechnology,RufusScaling}. Reactive came
first at the scale that matters most.

The central claim of this thesis is that a hybrid architecture
combining classical recommender algorithms with a reactive,
grounded LLM advisor can answer real shopping questions without
inheriting the trust failures of production systems, and that
every model and prompt choice in such an architecture can be
backed by measured evidence rather than intuition.

\section{Objectives}
\label{sec:intro-objectives}

This thesis sets out to design, build, and evaluate a hybrid
LLM-augmented shopping recommender with grounded output that
puts the central claim above into practice.

\begin{enumerate}
  \item Design a hybrid recommendation architecture in which a
    classical layer (Jaccard co-occurrence for ``frequently
      bought together'', weighted attribute similarity for
      ``similar products'', exponential time decay for trending,
    and a switching hybrid for personalized slots) generates
    candidates under a tight latency budget, while the LLM runs
    only on candidate sets, never on the raw catalog.

  \item Implement a reactive-first interaction model driven by
    lightweight client-side behavior signals (dwell, scroll,
    viewport, review engagement). Signals prepare context and
    surface a subtle availability indicator. They do not trigger
    LLM calls.

  \item Build a review summarization flow that surfaces
    conflicting buyer opinions with counts, fixing the failure
    most criticized in comparable production systems.

  \item Build an AI-assisted comparison flow that produces a
    structured tradeoff between two or three same-category
    products instead of a single ``best product'' verdict.

  \item Enforce hydration over generation as a structural defense
    against hallucination: the LLM emits product identifiers, the
    client renders live data from the catalog, and a validation
    layer checks every cited identifier, strips factual claims
    from free text, and verifies review-summary themes against
    the source corpus.

  \item Evaluate the system with a custom LLM evaluation harness
    that backs every configuration choice (model, reasoning
    effort, step budget, prompt variant) with measured evidence
    on grounded quality, latency, and cost.
\end{enumerate}

\section{Thesis Structure}
\label{sec:intro-structure}

The rest of the document is four chapters.

Chapter~\ref{chap:background} surveys the literature: classical
recommender systems, LLM-augmented recommendation, research on
human-AI trust and proactive versus reactive assistance, behavior
signals as proxies for intent, and a focused Amazon Rufus case
study that grounds the specific failure modes this system
addresses.

Chapter~\ref{chap:foundations} is the theory chapter. It defines
the building blocks the prototype uses: similarity measures
(Jaccard, weighted attribute scoring, cosine, exponential time
decay), the two-stage retrieve-and-rerank pattern, the LLM-agent
loop with tool use and structured output, the
hydration-over-generation principle as a formal defense against
hallucination, and the evaluation metrics used to compare
configurations.

Chapter~\ref{chap:system} documents the prototype. It opens with
the requirements and the overall architecture, then walks
through the data layer, the classical recommendation engine, the
behavior tracking and readiness pipeline, the LLM advisor with
its tools and prompt, the review summarization batch flow, the
AI-assisted comparison surface, the output validation pipeline,
and the frontend. The chapter closes with the evaluation harness
and the measured results across a nine-configuration sweep,
including pareto fronts on cost-quality and latency-quality and
a composite ranking that justifies the shipped configuration.

Chapter~\ref{chap:conclusions} summarizes what the prototype
demonstrates, what the measured evaluation supports, and the
limitations and directions for future work.
